% !TEX program = lualatex
\documentclass{ltjsarticle}
\usepackage{ketpic2e,ketlayer2e,graphicx,xcolor}
\usepackage{bm}


\usepackage{amsmath,amssymb}% amsfontsを読みに行くamssymb などを指定

% 通常の \mathbb を保持し，bboldフォントを使う
\DeclareMathAlphabet{\bbold}{U}{bbold}{m}{n}

% 自作フォント，CMMB10kubo.otfを\kuboldコマンドで使う
\usepackage{luatexja-fontspec}% \usepackage{fontspec}でもよい
\newfontfamily\CMMBkuboFont{CMMB10kubo.otf}
\DeclareRobustCommand{\kubold}[1]{\text{\CMMBkuboFont #1}}

% 数字にstixbbフォントを使う
\DeclareFontEncoding{LS1}{}{}
\DeclareFontSubstitution{LS1}{stix}{m}{n}
\DeclareFontFamily{LS1}{stixbb}{\skewchar\font127}
\DeclareFontShape{LS1}{stixbb}{m}{n}{<-> stix-mathbb}{}
\DeclareMathAlphabet{\stixbb}{LS1}{stixbb}{m}{n}

%\setmargin{15}{15}{20}{20}%% ketpicにより調整

% 用紙サイズと余白の設定
\usepackage[
  paperwidth=210mm,
  paperheight=165mm,
  left=15mm,
  right=15mm,
  top=5mm,
  bottom=5mm
]{geometry}

\pagestyle{empty}

\begin{document}
\Large\noindent
{\huge 【黒板文字を使う】}\\
黒板太字の小文字はベクトルを太字で書くときに使う。板書のときに使う簡略形を{\LaTeX}文書に使いたい。板書のときの簡略形は統一ルールがないように見える。ここで紹介するのは，一例とする。(これは，Lua\LaTeX でコンパイルしている。)

\textgt{参考：}amsfontsの\verb|\mathbb|による大文字\\
$\mathbb{ABCDEFGHIJKLMNOPQRSTUVWXYZ}$\\

\textgt{提案１：}大文字，小文字，数字をすべて bboldフォントで表示する。\par

\verb|\bbold| の確認（\verb|\bbold{a}|, \verb|\bbold{b}|, \verb|\bbold{c}|, ...）\par
大文字：$\bbold{ABCDEFGHIJKLMNOPQRSTUVWXYZ}$\par
小文字：$\bbold{abcdefghijklmnopqrstuvwxyz}$\qquad
数字：$\bbold{0123456789}$\\


\textgt{提案2：}大文字は amsfonts（\verb|\mathbb|），数字は stixフォント（\verb|\stixbb|）を使う。小文字は 久保が作成したフォント（\verb|CMMB10kubo.otf|）を使う。\par
大文字：$\mathbb{ABCDEFGHIJKLMNOPQRSTUVWXYZ}$（\verb|\mathbb|）\par
小文字：$\kubold{abcdefghijklmnopqrstuvwxyz}$\par
数字：$\stixbb{0123456789}$（\verb|\stixbb|）\\
※ cmmib10.pfb を参考に，久保が \verb|CMMB10kubo.otf| を作成した。
このフォントは，ワードなどtex以外のアプリでも使える。



\vspace{\baselineskip}

※ 提案1, 2のプリアンブルは，texソースを参照。\par

\end{document}